% Source: https://artofproblemsolving.com/community/c3h2650199_this_probably_already_exists_but_whatever
% Stable identity: aops-topic-2650199-q01
\textbf{(a)} A result is a pair $(e,c)$ of nonnegative integers ($e$ blanks, $c$ correct) with $e+c\le P$; there are $\dbinom{P+2}{2}$ such pairs. The score is $ea+cb$, and
\[ea+cb=e'a+c'b\iff\left(e-e'\right)a=\left(c'-c\right)b\iff e-e'=tb,\ c'-c=ta\text{ for some integer }t,\]
since $\gcd(a,b)=1$. So pairs with equal scores form chains $(e+tb,\,c-ta)$. The conditions for such a pair to be a valid result are linear inequalities in $t$, so each chain is a run of consecutive values of $t$. Moving along a chain changes $e+c$ by $t(b-a)$.

Suppose $b\ge a$, so $x=b$. (If $a=b$ then $a=b=1$, and the argument below still works.) In each chain exactly one pair, the one with the largest $t$, has no successor $(e+b,c-a)$ in the valid region. So the number of distinct scores equals the number of pairs \emph{without} a valid successor. A pair has a valid successor exactly when $c\ge a$ and $(e+b)+(c-a)\le P$. Writing $c=a+c''$, these pairs correspond to the pairs $(e,c'')$ with $e+c''\le P-b$, of which there are $\dbinom{P+2-b}{2}$ (this is $0$ when $P+2-b\in\{0,1\}$). Hence the number of scores is $\dbinom{P+2}{2}-\dbinom{P+2-x}{2}$. If $a>b$, the same argument with successor $(e-b,c+a)$ gives $\dbinom{P+2-a}{2}$ pairs, so the formula is the same.

Finally,
\[\binom{P+2}{2}-\binom{P+2-x}{2}=\dfrac{(P+2)(P+1)-(P+2-x)(P+1-x)}{2}=\dfrac{x(2P+3)-x^2}{2}.\]
(The assumption $x\le P+2$ is needed: for $P=1$, $a=1$, $b=4$ the scores are $0,1,4$, but the formula gives $2$.) $\blacksquare$

\textbf{(b)} Scores are $1.5(e+4c)$, so this is the case $a=1$, $b=4$, $P=25$: there are $\dfrac{4\cdot53-16}{2}=\boxed{98}$ possible scores.
